\section{KMP Algorithm and the LPS Array}
\label{sec:str-kmp-lps}

\problemstatement{Compute the longest-proper-prefix-which-is-also-suffix
(LPS/prefix-function) array of a pattern, and use it to find all occurrences
of the pattern in a text in linear time.}

\begin{patternbox}
\textbf{KMP} avoids re-comparing characters after a mismatch by precomputing
the \textbf{LPS array}: for each pattern prefix, the length of its longest
border. On a mismatch, the pattern shifts so the matched border aligns,
never moving the text pointer backward.
\end{patternbox}

\subsection*{Borders let you resume, not restart}

The LPS array $\pi[i]$ is the longest proper prefix of
$\textit{pattern}[0..i]$ that is also a suffix. Build it by extending or
falling back along previously computed borders. During matching, when a
mismatch occurs after matching $j$ characters, instead of restarting you set
$j=\pi[j-1]$ -- the next-longest border already matched -- and try again. The
text pointer only moves forward, giving $O(n+m)$.

\begin{ideabox}[The LPS Says Where to Resume After a Mismatch]
A border of the matched portion is guaranteed to align with the text, so
after a mismatch KMP jumps the pattern to that border rather than
re-examining text characters. Precomputing all borders makes each jump
$O(1)$ amortised.
\end{ideabox}

\subsection*{Algorithm}

\begin{enumerate}
  \item Build $\pi$ for the pattern: extend on match, fall back via
        $\pi[j-1]$ on mismatch.
  \item Scan the text with a pointer $j$ into the pattern; on match advance
        both, on mismatch set $j=\pi[j-1]$.
  \item When $j$ reaches $|\textit{pattern}|$, record an occurrence and set
        $j=\pi[j-1]$.
\end{enumerate}

\subsection*{Dry run}

\dryrun{pattern \texttt{"aba"}, text \texttt{"ababa"}.}

\begin{itemize}
  \item $\pi=[0,0,1]$. Matching finds \texttt{"aba"} at index $0$; after a
        match, fall back to reuse the overlap and match again at index $2$.
  \item Occurrences at $0$ and $2$.
\end{itemize}

\subsection*{C++ implementation}

\begin{lstlisting}
#include <bits/stdc++.h>
using namespace std;

vector<int> buildLPS(const string& p) {
    int m = p.size();
    vector<int> pi(m, 0);
    for (int i = 1; i < m; i++) {
        int j = pi[i - 1];
        while (j > 0 && p[i] != p[j]) j = pi[j - 1];   // fall back on mismatch
        if (p[i] == p[j]) j++;
        pi[i] = j;
    }
    return pi;
}

vector<int> kmp(const string& text, const string& pat) {
    vector<int> pi = buildLPS(pat), res;
    int n = text.size(), m = pat.size(), j = 0;
    for (int i = 0; i < n; i++) {
        while (j > 0 && text[i] != pat[j]) j = pi[j - 1];
        if (text[i] == pat[j]) j++;
        if (j == m) { res.push_back(i - m + 1); j = pi[j - 1]; }
    }
    return res;
}

int main() {
    for (int p : kmp("ababa", "aba")) cout << p << " ";
    cout << "\n";                                  // 0 2
    return 0;
}
\end{lstlisting}

\begin{complexitybox}
\textbf{Time Complexity.} $O(n+m)$ -- LPS build and matching are each
linear. \textbf{Space Complexity.} $O(m)$ for the LPS array.
\end{complexitybox}

\begin{takeawaybox}
KMP precomputes borders (the LPS array) so a mismatch resumes from the
longest overlap instead of restarting, keeping the text pointer moving
forward. The LPS array itself answers border/period questions and powers the
next two problems.
\end{takeawaybox}
